\section{Introduction}
\label{sec:intro}
% Base: Ivan Ezhov's version of 13 Nov 2026 (Overleaf), with Paula's term questions answered in place.
Inpainting asks a model to fill a hidden region of an image from its visible surroundings.
In medical imaging it replaces content that is missing or must not be used: for brain tumor patients, even the first scan shows the lesion, and inpainting healthy-looking tissue into it lets tools built for healthy brains, such as tissue segmentation and brain extraction, run on the scan~\cite{brats2023}.
Such a fill should show this patient's anatomy, but an unsure model can reduce its error by producing blurry, population-average content instead.
The fills are scored by their distance from the original content, typically using metrics such as \ac{MSE}, \ac{SSIM}, or \ac{LPIPS}~\cite{zhang2018lpips}.
A low distance can mean that the fill recovered detail that belongs to this individual, or only that it is a generic, population-average fill that is close on average.
A fill can also look plausible yet show someone else's face, eye, or brain.
A single score cannot distinguish these cases.
For a model output, the score a metric should give is unknown, so the metric cannot be checked on it.
We therefore audit the metrics themselves.
We instead score fills whose relationship to the target is known by design: real images, averages of real images, and controls built from them.
The audit is deterministic and fully data-driven, and what a metric rewards here does not depend on any particular model.

Human anatomy provides such references.
A second image of the same person, such as a second \ac{MRI} session or another photograph, preserves the person's anatomy but might change the acquisition parameters.
The opposite side of the same body, such as the other eye, hemisphere, or half of the face, keeps the person but shows a mirrored structure.
Genetic relatives may share some anatomy; unrelated people share only what is common to the population; and averaging across many unrelated people preserves the common structure without individual variation.

We arrange such references into the \emph{identity ladder}, a calibration scale for inpainting metrics; each kind of reference is a \emph{rung}, and its expected position is known before any metric is computed.
We call the difference between a person's image and the dataset average their \emph{individual detail}.
Three constructed controls separate losing detail from getting it wrong: the \emph{softened same person} is the person's own second image with its individual detail scaled down toward the average; \emph{wrong detail} adds a random person's individual detail of the target's strength to the average; and \emph{noise} adds white noise of that strength.
Finally, unrelated people chosen solely because their visible surroundings resemble the target's show what a fill can achieve without knowing the person.
We ask four research questions.
\textbf{RQ1:} Does a metric place the same person closer than a random person?
\textbf{RQ2:} Does it penalize partial averaging toward the dataset average?
\textbf{RQ3:} Does it separate lost detail (the right person, made blurry) from wrong detail (sharp detail from someone else)?
\textbf{RQ4:} How close can unrelated people chosen from the visible surroundings alone come to the same person?

Our contributions are the identity ladder with its three control conditions, an audit of ten metric variants from four families on nine image sets from eight cohorts of brain \ac{MRI}, retinal photographs, faces, and irises, and the positions of four pretrained inpainting models on that ladder.
No tested metric passes both RQ1 and RQ2: metrics computed on the filled region reward partial averaging, and set-level metrics, \ac{FID}~\cite{heusel2017fid} and \ac{FWD}~\cite{fwd2025}, barely separate the right person from a random person.
That the conditional mean minimizes expected squared error is known~\cite{blau2018}; the audit adds the size of this reward on real anatomy, the opposite failure of set-level metrics, a retrieval baseline built from the context alone whose distance depends on the region size, and directional tests in each image set.
Inpainting algorithm benchmarks, such as the \ac{BraTS} brain inpainting challenge~\cite{brats2023}, should report these references and controls next to their own metrics.
Inpainting is also a convenient test bed for conditional generation: the condition is fixed, and the content the model should produce exists, so whatever a metric rewards here, it rewards wherever generated content is scored against a single reference.
